\documentclass[12pt,a4paper]{article}

\usepackage[margin=1in]{geometry}
\usepackage[T1]{fontenc}
\usepackage[utf8]{inputenc}
\usepackage{lmodern}
\usepackage{graphicx}
\usepackage{booktabs}
\usepackage{tabularx}
\usepackage{longtable}
\usepackage{array}
\usepackage{hyperref}
\usepackage{xcolor}
\usepackage{caption}
\usepackage{float}

\hypersetup{
  colorlinks=true,
  linkcolor=blue!60!black,
  urlcolor=blue!60!black,
  citecolor=blue!60!black,
  pdftitle={Final Project Report - STEM Learning Platform with GenAI},
  pdfauthor={Project Team}
}

\graphicspath{
  {../docs/mermaid-png/root/}
  {../docs/mermaid-png/web/}
  {../docs/mermaid-png/backend/}
  {../docs/mermaid-png/supabase/}
}

\setlength{\parindent}{0pt}
\setlength{\parskip}{0.65em}
\renewcommand{\arraystretch}{1.2}

\newcommand{\ProjectTitle}{STEM Learning Platform with GenAI}
\newcommand{\ReportTitle}{Final Project Report}
\newcommand{\CourseCode}{COMP3122}
\newcommand{\PreparedBy}{Project Team}
\newcommand{\ReportDate}{April 2026}

\begin{document}
\pagenumbering{gobble}

\begin{titlepage}
\centering
\vspace*{2cm}
{\Large \CourseCode \par}
\vspace{1.2cm}
{\Huge \bfseries \ReportTitle \par}
\vspace{0.8cm}
{\LARGE \ProjectTitle \par}
\vspace{1.5cm}
\begin{tabular}{>{\bfseries}p{4cm}p{8cm}}
Prepared by & \PreparedBy \\
Date & \ReportDate \\
Repo snapshot & 2026-04-05 \\
\end{tabular}
\vfill
{\large
This report documents the current implemented state of the repository, including the
web application, Python backend, Supabase schema and Edge Functions, testing strategy,
deployment topology, and role-based user workflows.
\par}
\end{titlepage}

\clearpage
\pagenumbering{arabic}
\tableofcontents
\newpage

\section{Executive Summary}

This project delivers a production-oriented STEM learning platform that integrates
Generative AI into classroom workflows rather than treating AI as a detached chatbot.
Its main idea is to let instructors upload teaching materials, convert those materials
into a structured \emph{Course Blueprint}, and use that blueprint as the source of truth
for downstream student learning activities such as grounded chat, quizzes, flashcards,
assignment review, and class-level insights.

The implemented system is split into three cooperating runtime surfaces:
\begin{itemize}
  \item a \textbf{Next.js 16} web application for role-aware user experience, route handling,
  and server actions;
  \item a \textbf{Python FastAPI} backend for AI orchestration, class workflows, analytics,
  and chat memory management; and
  \item a \textbf{Supabase} project for authentication, PostgreSQL persistence, row-level
  security, storage, queue-backed material processing, and guest sandbox lifecycle support.
\end{itemize}

From a product perspective, the platform supports:
\begin{itemize}
  \item \textbf{instructors / graders} who create classes, upload materials, generate and
  publish blueprints, create assignments, review submissions, and inspect analytics;
  \item \textbf{students} who join classes with join codes, use the published blueprint,
  complete assignments, and learn through grounded AI chat; and
  \item \textbf{guest users} who can explore the product through a realistic sandboxed demo
  flow backed by Supabase Anonymous Auth and sandbox-cloned data.
\end{itemize}

The repository also demonstrates a meaningful software engineering process. The current
snapshot contains 401 commits, 92 merge commits, 23 active database migrations, a multi-layer
test strategy, and a documented deployment topology across web, backend, and Supabase services.
The result is not a toy prototype, but a coordinated educational platform with clear subsystem
boundaries and operational considerations.

\section{Project Background and Motivation}

\subsection{Course and problem context}

The supplied project overview and background materials position the project in the context of
\emph{Education 4.0}, AI-empowered personalized learning, and the practical pain points faced by
teachers. The brief explicitly asks teams to avoid building ``just another AI Q\&A app'' and
instead create a platform with meaningful educational value, differentiated features, and a
well-considered software engineering process.

The background deck highlights several recurring teaching needs:
\begin{itemize}
  \item reducing instructor workload in lesson preparation and grading;
  \item generating learning activities such as quizzes and guided exercises;
  \item using AI to support more personalized and interactive learning;
  \item grounding AI output in trustworthy course content rather than generic internet-style chat.
\end{itemize}

These needs strongly influenced the design of this system. Rather than exposing a general-purpose
chatbot, the platform is built around a teacher-controlled workflow where uploaded course materials
are curated into a blueprint before AI-generated learning activities are released to students.

\subsection{Project goals}

Based on the repository documentation and implementation, the core goals of the project are:
\begin{itemize}
  \item transform raw teaching materials into structured, reusable course knowledge;
  \item keep instructors in editorial control of AI outputs;
  \item provide students with guided, class-specific learning experiences;
  \item support realistic classroom operations such as assignment creation, review, and feedback;
  \item demonstrate a robust engineering process including testing, deployment, and iteration.
\end{itemize}

\subsection{What makes this project distinctive}

Several aspects make this project more substantial than a basic GenAI demo:
\begin{itemize}
  \item the \textbf{Course Blueprint lifecycle} introduces an approval and publishing checkpoint;
  \item \textbf{all AI generation routes through a dedicated Python service}, not ad hoc frontend calls;
  \item \textbf{material processing is asynchronous and queue-backed}, which is closer to production reality;
  \item \textbf{teacher intelligence features} extend beyond content generation into analytics and adaptive briefs;
  \item \textbf{guest mode} is implemented as a true sandboxed experience rather than a fake slideshow demo.
\end{itemize}

\section{System Overview and Implemented Features}

\subsection{High-level system picture}

Figure~\ref{fig:system-overview} shows the top-level structure of the current platform.

\begin{figure}[H]
  \centering
  \includegraphics[width=\textwidth]{README-diagram-01.png}
  \caption{Current high-level platform architecture.}
  \label{fig:system-overview}
\end{figure}

At a high level, teacher, student, and guest users all enter through the web application.
The web layer handles UI rendering and user flow orchestration, while the backend is responsible
for AI-heavy logic and the Supabase layer provides persistence, storage, security, and async jobs.

\subsection{Primary user groups}

\begin{table}[H]
\centering
\caption{Primary user groups in the current release}
\begin{tabularx}{\textwidth}{>{\bfseries}p{3cm}X}
\toprule
Role & Implemented responsibilities \\
\midrule
Instructor / Grader & Create classes, upload materials, generate and curate blueprints, create assignments, review student work, monitor chat activity, inspect class insights, use the adaptive teaching brief, and preview the student experience. \\
Student & Join classes using a join code, access the published blueprint, use always-on AI chat, complete chat/quiz/flashcards assignments, and track progress from the student dashboard. \\
Guest & Enter a sandbox demo, explore both teacher and student views through role switching, and safely evaluate the product without touching real user data. \\
\bottomrule
\end{tabularx}
\end{table}

\subsection{Feature set}

\begin{table}[H]
\centering
\caption{Major implemented product features}
\begin{tabularx}{\textwidth}{>{\bfseries}p{3.3cm}X}
\toprule
Feature area & Current implementation \\
\midrule
Authentication & Email/password sign-in for permanent accounts, email verification, password reset, and guest access via Supabase Anonymous Auth. \\
Class management & Teacher class creation, student join-by-code, role-aware dashboards, and class-scoped navigation. \\
Materials workflow & PDF, DOCX, and PPTX upload; asynchronous extraction/chunking/embedding; preview, download, and delete actions. \\
Blueprint workflow & Draft, overview, and published blueprint lifecycle with teacher editing and publish control. \\
Learning activities & Chat assignments, quiz assignments, flashcard assignments, plus open-practice class chat. \\
Teacher review & Assignment review screens for chat, quiz, and flashcards, including feedback-saving flows. \\
Analytics & Class intelligence dashboard with summaries, charts, student drill-downs, and intervention suggestions. \\
Teaching support & Adaptive teaching brief widget summarizing what needs attention and recommended next actions. \\
Visual explanation & Generative canvas support for charts and diagram-style responses in chat and analytics contexts. \\
Guest mode & Sandbox provisioning, role switching, reset flow, quota management, expiry, and cleanup. \\
\bottomrule
\end{tabularx}
\end{table}

\subsection{Blueprint-centered value flow}

The most important end-to-end educational flow is shown in Figure~\ref{fig:value-flow}.

\begin{figure}[H]
  \centering
  \includegraphics[width=\textwidth]{README-diagram-02.png}
  \caption{Blueprint-centered value flow from teacher materials to student learning and review.}
  \label{fig:value-flow}
\end{figure}

This flow is central to the platform's pedagogy:
\begin{enumerate}
  \item the instructor uploads materials;
  \item the system processes and embeds those materials;
  \item the AI generates a structured blueprint;
  \item the instructor edits and approves the blueprint;
  \item the published blueprint becomes the basis for activities;
  \item students learn through grounded interactions;
  \item instructors review outputs and analyze class performance.
\end{enumerate}

\section{System Design}

\subsection{Technology stack}

\begin{table}[H]
\centering
\caption{Core implementation stack}
\begin{tabularx}{\textwidth}{>{\bfseries}p{3.5cm}X}
\toprule
Layer & Technology \\
\midrule
Frontend & Next.js 16 App Router, React 19, TypeScript, Tailwind CSS 4, shared UI primitives, Motion for React, Lucide icons \\
Backend & Python FastAPI, httpx, structured Pydantic schemas, provider adapters \\
Database and auth & Supabase PostgreSQL, Supabase Auth, Row Level Security \\
Storage and async processing & Supabase Storage, pgmq queue, pg\_cron, Supabase Edge Functions \\
Testing & Vitest, React Testing Library, Python unittest, Playwright \\
Deployment model & Separate web and backend deployments plus hosted Supabase \\
\bottomrule
\end{tabularx}
\end{table}

\subsection{Architectural decomposition}

The current repository is intentionally partitioned into four main top-level areas:
\begin{itemize}
  \item \texttt{web/}: user-facing routes, UI components, middleware, and server actions;
  \item \texttt{backend/}: AI orchestration, chat workspace logic, class workflows, analytics;
  \item \texttt{supabase/}: SQL migrations, local config, and Edge Functions;
  \item \texttt{tests/}: Playwright end-to-end coverage running against a deployed target URL.
\end{itemize}

This separation matters because different responsibilities are intentionally placed in the layer
best suited for them. The web app owns rendering and navigation; the Python backend owns AI and
workflow-heavy service logic; Supabase owns security, persistence, storage, and background jobs.

\subsection{Web application design}

The web application is not just a collection of pages. It implements a role-aware interface model
with separate teacher and student dashboards, a persistent class shell, and route-level auth
protection.

Important web-layer characteristics include:
\begin{itemize}
  \item role-specific dashboards at \texttt{/teacher/dashboard} and \texttt{/student/dashboard};
  \item unified auth via a shared \texttt{AuthSurface} used in modal and dedicated-page forms;
  \item middleware-enforced redirects for unauthenticated, unverified, and expired guest sessions;
  \item server actions for class creation, join-by-code, materials, assignments, and settings;
  \item a shared design system with centralized icons, tokens, and motion presets.
\end{itemize}

On the teacher side, the class overview acts as a control hub for materials, assignments,
teaching brief, and monitoring. On the student side, the class workspace is organized around
focus widgets for AI chat, assignments, and the published blueprint. A teacher can also preview
the student experience without changing accounts.

\subsection{Python backend design}

The backend is the sole AI orchestration boundary in the current implementation. This is an
important design choice because it avoids scattering prompt construction, provider credentials,
and response normalization across the frontend.

The backend currently provides endpoints and domain logic for:
\begin{itemize}
  \item blueprint generation;
  \item quiz generation;
  \item flashcard generation;
  \item grounded chat;
  \item chat workspace sessions and message persistence;
  \item visual canvas generation;
  \item class create and class join;
  \item materials dispatch;
  \item analytics, teaching brief, and natural-language data queries.
\end{itemize}

All normal backend responses follow a consistent envelope:
\begin{center}
\texttt{\{ ok, data, error, meta \}}
\end{center}

This design improves consistency at the integration boundary. The frontend adapters can handle
success and failure cases uniformly, while the backend retains freedom to evolve internal provider
logic and prompt design.

\subsection{AI orchestration design}

The repository currently supports OpenRouter, OpenAI, and Gemini as pluggable providers.
Provider choice is environment-driven, and the backend can fall back across providers when needed.

This design has several advantages:
\begin{itemize}
  \item the web app remains provider-agnostic;
  \item deployment environments can switch provider preference through configuration;
  \item provider failures can be isolated inside the backend;
  \item request deadlines are enforced centrally rather than inconsistently per route.
\end{itemize}

The current AI domains are broader than content generation alone. In addition to blueprint, quiz,
and flashcard generation, the system also supports analytics narratives, teaching brief synthesis,
chat memory handling, and visual canvas specifications.

\subsection{Chat workspace and memory design}

One of the most technically interesting parts of the system is the chat workspace. It is not
implemented as a purely stateless prompt-response loop. Instead, the backend includes a dedicated
chat workspace subsystem with participant listing, session creation and archive, paginated message
history, and context compaction for long-running conversations.

\begin{figure}[H]
  \centering
  \includegraphics[width=\textwidth]{ARCHITECTURE-diagram-08.png}
  \caption{Chat workspace flow with grounded context retrieval and persisted assistant responses.}
  \label{fig:chat-workspace}
\end{figure}

This design is significant for two reasons. First, it better reflects real classroom use, where
students may return to a conversation over time rather than ask only one question. Second, it
forces the system to deal with long-context management, not just one-shot prompt construction.

\subsection{Material processing pipeline}

Material ingestion is intentionally asynchronous. Uploading a file creates a material record and
stores the file, but extraction, chunking, and embedding happen through a queue-backed worker path
instead of blocking the user's request.

\begin{figure}[H]
  \centering
  \includegraphics[width=\textwidth]{DEPLOYMENT-diagram-02.png}
  \caption{Asynchronous material-processing pipeline.}
  \label{fig:material-pipeline}
\end{figure}

This was a good architectural decision for a course project because it addresses a real system
concern: extraction and embedding are too slow and failure-prone to treat as a trivial synchronous
request. The queue-backed design allows better retry behavior, clearer UI states (\texttt{processing},
\texttt{ready}, \texttt{failed}), and a cleaner separation between user interaction and background work.

\subsection{Analytics and adaptive teaching support}

The platform includes teacher-facing intelligence features, which meaningfully extend it beyond
content generation.

The class intelligence dashboard includes:
\begin{itemize}
  \item summary cards for student count, average score, completion rate, and risk indicators;
  \item topic-level performance charts;
  \item Bloom-level breakdown visualization;
  \item student engagement views and student overview tables;
  \item AI-generated narrative interpretation and intervention suggestions;
  \item natural-language data query support for generating chart specifications.
\end{itemize}

The adaptive teaching brief is a complementary feature. Rather than exposing a large dashboard
all the time, it provides a concise memo-style summary directly on the class overview page, with
recommended next steps, likely misconceptions, and follow-up activity suggestions.

\subsection{Security and access-control design}

Security is handled across all three runtime surfaces:
\begin{itemize}
  \item \textbf{web layer}: route protection, email verification enforcement, guest route confinement;
  \item \textbf{backend layer}: service authentication, user bearer-token verification, guest AI guardrails;
  \item \textbf{Supabase layer}: row-level security, teacher/enrollment policies, and sandbox-aware isolation.
\end{itemize}

The current implementation documents several security-conscious choices:
\begin{itemize}
  \item permanent users are limited to teacher or student account types;
  \item \texttt{profiles.account\_type} is immutable after signup;
  \item guest sessions are scoped by \texttt{sandbox\_id};
  \item cursor inputs to PostgREST are validated before interpolation;
  \item service-to-service auth between web and backend uses a shared backend API key.
\end{itemize}

\subsection{Guest sandbox design}

Guest mode is one of the strongest differentiators in the project. Instead of presenting a static
marketing demo, the system provisions a real anonymous session and clones a seeded classroom into
a sandboxed data space. Guests can switch between teacher and student views, reset the demo, and
observe realistic product behavior without accessing real user data.

The current quota model enforces:
\begin{itemize}
  \item a global cap on active guest sessions;
  \item a per-hour cap on newly created guest sessions;
  \item backend limits on guest AI usage and concurrent in-flight AI requests;
  \item expiry after 32 hours maximum or 8 hours of inactivity.
\end{itemize}

This is a thoughtful design because it balances demo accessibility with safety, resource control,
and cleanup reliability.

\section{Software Development Process}

\subsection{Requirements and scope shaping}

The initial project brief encouraged teams to investigate stakeholder needs, review existing
platforms, and avoid shallow AI-first implementations. The current repository reflects that
direction. The system is not scoped as a universal learning platform, but as a structured STEM
classroom assistant centered on teacher materials, editorial control, and student-facing learning
activities.

A practical strength of the project is its restraint. The team did not attempt to implement every
possible educational feature. Instead, it converged on a coherent workflow:
\begin{itemize}
  \item content ingestion;
  \item blueprint generation and publishing;
  \item student chat and assignments;
  \item review and analytics;
  \item guest evaluation flow.
\end{itemize}

This focused scope likely helped the team build a more integrated product instead of a fragmented
set of unrelated AI experiments.

\subsection{Design process}

The design process visible in the repository appears iterative and documentation-driven.
The project contains:
\begin{itemize}
  \item top-level design, architecture, UI/UX, and deployment documents;
  \item internal reference notes under \texttt{.codex/};
  \item Mermaid diagrams exported to PNG for communication and reporting;
  \item commit history showing multiple phases of refinement rather than one-shot implementation.
\end{itemize}

The UI/UX work also appears to have been treated as a first-class concern. The repository contains
significant evidence of design-token work, navigation restructuring, auth UX redesign, teacher/student
flow differentiation, and consistency hardening. This is valuable because educational software quality
depends not only on AI correctness, but also on clarity, trust, and ease of use.

\subsection{Implementation phases observed in git history}

The commit history provides a useful implementation timeline. Table~\ref{tab:timeline} summarizes
the major development phases visible from the repository.

\begin{table}[H]
\centering
\caption{Major development phases observed from git history}
\label{tab:timeline}
\begin{tabularx}{\textwidth}{>{\bfseries}p{2.8cm}X}
\toprule
Period & Main progression visible in the repository \\
\midrule
2026-02-03 to 2026-02-05 & Repository setup, initial schema and RLS, core class flows, materials pipeline, AI adapters, blueprint generation, and first UI scaffolding. \\
2026-02-06 to 2026-02-10 & Vertical slices for AI chat, quiz generation, flashcards, auth hardening, and always-on class chat workspace. \\
2026-02-22 to 2026-03-02 & Material worker stabilization, queue-backed processing, dashboard and class-shell UX improvements, theme/token refactoring, and stronger unit-test coverage. \\
2026-03-15 to 2026-03-24 & Teacher class intelligence dashboard, generative canvas, teaching brief backend and widget, content review improvements, and more polished materials preview behavior. \\
2026-03-26 to 2026-03-28 & Guest mode foundation and hardening, E2E expansion, shared auth redesign, email template work, OTP-based password change flow, and documentation refresh. \\
2026-04-04 to 2026-04-05 & Guest quota redesign, concurrency improvements, final bug fixing, environment sync, and diagram/documentation alignment with runtime reality. \\
\bottomrule
\end{tabularx}
\end{table}

This sequence suggests a healthy development pattern:
\begin{enumerate}
  \item build the skeleton;
  \item deliver core vertical slices;
  \item strengthen architecture and UX;
  \item add analytics and guest mode;
  \item harden the system and documentation near submission.
\end{enumerate}

\subsection{Collaboration workflow}

The collaboration model that is directly observable from the repository is based on GitHub-style
branching, pull requests, review, and merge integration.

Evidence visible in the current snapshot includes:
\begin{itemize}
  \item 401 commits in total;
  \item 92 merge commits;
  \item repeated feature-branch merge patterns such as guest mode, UI improvements, analytics,
  teaching brief, and auth flows;
  \item regular documentation updates that track design decisions, rollout notes, and lessons learned.
\end{itemize}

This indicates an iterative workflow rather than a single large code dump. The team appears to
have used:
\begin{itemize}
  \item feature branches for major workstreams;
  \item pull requests to integrate features and fixes;
  \item tests and documentation updates alongside implementation changes;
  \item structured plans and handoff notes for ongoing work.
\end{itemize}

At the same time, one accuracy caveat is important: the visible git author history is concentrated
under a small number of identities, especially one primary contributor identity. For that reason,
this report documents the \emph{repository-observable collaboration process} rather than making
unsupported claims about each team member's exact offline contribution split. If the course requires
a per-member effort table, that should be appended using the team's own meeting notes or task records.

\subsection{Use of GenAI during development}

This project is itself a GenAI platform, but the repository also suggests that GenAI-assisted
development was part of the engineering workflow. The documentation explicitly discusses Codex- and
agent-assisted workflows, planning notes, and iterative bug fixing. From a software engineering
perspective, this is useful to reflect on because GenAI appears to have been used as:
\begin{itemize}
  \item a coding accelerator for scaffolding and incremental implementation;
  \item a design aid for exploring architecture and UI alternatives;
  \item a debugging assistant for reproducing and fixing workflow bugs;
  \item a documentation aid for maintaining architecture and usage guides.
\end{itemize}

However, the repository also shows good evidence of human judgment:
\begin{itemize}
  \item repeated refactors and bug-fix passes after initial implementations;
  \item hardening around access control, redirect behavior, and quota logic;
  \item correction of real production-style issues such as signed URL previewing, timeout handling,
  guest cleanup, and environment misconfiguration;
  \item a sustained effort to align documentation with the actual runtime state.
\end{itemize}

This combination is exactly the kind of responsible GenAI-assisted development that a modern course
project should demonstrate: AI helps with speed, but humans still own design quality, correctness,
and final judgment.

\section{Testing, Quality Assurance, and Deployment}

\subsection{Testing strategy}

The repository implements a layered testing strategy rather than relying on only one test type.

\begin{table}[H]
\centering
\caption{Current testing layers in the repository}
\begin{tabularx}{\textwidth}{>{\bfseries}p{3.3cm}X>{\raggedleft\arraybackslash}p{2.2cm}}
\toprule
Layer & What it covers & Evidence in repo \\
\midrule
Web unit and integration tests & Server actions, auth flows, route guards, material actions, blueprint logic, chat flows, guest helpers, UI primitives, validation helpers & 60 test files under \texttt{web/src} \\
Backend unit tests & Providers, schemas, class flows, blueprint generation, quiz, flashcards, materials, analytics, chat workspace, guest rate limiting & 17 Python test files under \texttt{backend/tests} \\
End-to-end tests & Teacher navigation, student navigation, login, join class, settings, sign-out, and guest entry against a deployed environment & 9 Playwright specs under \texttt{tests/e2e} \\
\bottomrule
\end{tabularx}
\end{table}

Representative quality signals include:
\begin{itemize}
  \item guest middleware tests and guest sandbox tests;
  \item blueprint action tests;
  \item materials extraction and retrieval tests;
  \item chat generation, compaction, and workspace tests;
  \item analytics and teaching brief action tests;
  \item Playwright flows for both teacher and student navigation.
\end{itemize}

\subsection{Why the testing approach is appropriate}

This testing distribution is appropriate for the architecture:
\begin{itemize}
  \item unit tests catch domain logic regressions close to the code;
  \item web integration tests validate server actions and route behavior;
  \item backend tests verify AI-related orchestration and guardrails without needing full browser runs;
  \item Playwright tests confirm that the deployed user flows still work end to end.
\end{itemize}

For a project involving multiple services and stateful workflows, this layered approach is far more
credible than relying on only manual demonstration.

\subsection{Deployment topology}

The deployed system consists of a frontend deployment, a separate backend deployment, and a hosted
Supabase project. Figure~\ref{fig:deployment-topology} captures this structure.

\begin{figure}[H]
  \centering
  \includegraphics[width=\textwidth]{DEPLOYMENT-diagram-01.png}
  \caption{Hosted deployment topology across frontend, backend, and Supabase services.}
  \label{fig:deployment-topology}
\end{figure}

This split has practical benefits:
\begin{itemize}
  \item frontend issues can be distinguished from backend issues;
  \item AI orchestration can evolve independently from the UI deployment;
  \item Supabase-specific operational tasks such as migrations and Edge Function secrets remain explicit;
  \item the team can test preview or staging environments in a more controlled way.
\end{itemize}

\subsection{Deployment process}

The documented deployment workflow includes:
\begin{enumerate}
  \item configuring one Supabase project per environment;
  \item applying the active migration set;
  \item deploying the \texttt{material-worker} and \texttt{guest-sandbox-cleanup} Edge Functions;
  \item configuring Edge Function secrets and Vault-backed values;
  \item deploying the Python backend;
  \item deploying the web application;
  \item running post-deploy smoke tests for auth, teacher flows, student flows, insights, and guest mode.
\end{enumerate}

Operationally, this is a mature approach for a course project because it recognizes that deployment
is not just ``push frontend to Vercel''; the data plane and async workers also have to be correct.

\subsection{Database evolution}

The current migration history shows 23 active migrations, which is another sign that the project
evolved iteratively rather than remaining static after an initial schema dump.

Major migration themes include:
\begin{itemize}
  \item initial schema and RLS foundation;
  \item quiz integrity and auth hardening;
  \item canonical blueprint snapshots;
  \item always-on class chat and compaction;
  \item queue-backed material processing;
  \item class insights and teaching brief snapshots;
  \item guest mode schema, seed data, cleanup, and quota redesign.
\end{itemize}

This migration history is important because it mirrors the product's evolution in a traceable,
reviewable form.

\section{Lessons Learned and Reflection}

\subsection{Technical lessons}

The repository's maintained documentation highlights several concrete lessons that are worth
reflecting on.

\begin{itemize}
  \item \textbf{A clean architecture boundary matters.} Routing all AI through the Python backend
  keeps provider-specific complexity out of the UI and gives the system a stable application-facing
  contract.

  \item \textbf{Async processing is necessary for realistic content pipelines.} Material extraction
  and embedding do not belong in a naive synchronous upload request if the goal is a reliable system.

  \item \textbf{Guest mode is harder than it looks.} A believable demo mode requires session
  quotas, expiry, cleanup, sandbox isolation, and backend-side ownership checks. A demo that looks
  simple in the UI can still be architecturally complex.

  \item \textbf{Deployment secrets must live in the right place.} The project explicitly notes that
  Supabase Edge Functions cannot rely on Vercel secrets or Supabase Vault alone; function-level
  secrets must be configured in Supabase Edge Function Secrets.

  \item \textbf{Subtle browser and CSS issues can become product bugs.} Examples in the repository
  include iframe preview limitations caused by storage response headers and dialog-centering issues
  caused by Tailwind v4's use of the standalone CSS \texttt{translate} property.
\end{itemize}

\subsection{Process lessons}

From a software engineering process perspective, several lessons stand out:
\begin{itemize}
  \item keeping design, architecture, deployment, and UI docs current reduces confusion late in the project;
  \item documentation is especially valuable when a project spans frontend, backend, and database layers;
  \item a branch-and-merge workflow helps isolate major features such as guest mode and auth redesign;
  \item testing should expand together with complexity, not after the system is already unstable.
\end{itemize}

\subsection{What the project did well}

In its current form, the project demonstrates several strengths:
\begin{itemize}
  \item a coherent product concept instead of a scattered list of AI features;
  \item strong teacher-centered workflow design;
  \item real effort on UI quality and navigation clarity;
  \item a credible multi-service architecture;
  \item meaningful testing and deployment documentation.
\end{itemize}

\subsection{What could still be improved}

A professional reflection should also note realistic limitations:
\begin{itemize}
  \item the production UI currently has teacher/student roles but no dedicated admin console;
  \item the operational surface is sophisticated, so monitoring and observability could still grow further;
  \item the visible git author concentration suggests that a clearer contribution matrix may still be needed for academic assessment;
  \item broader activity types beyond chat, quiz, and flashcards remain future-facing rather than fully shipped.
\end{itemize}

\section{User Manual}

\subsection{Administrator / System Maintainer}

\textbf{Important accuracy note:} the current release does \emph{not} expose a dedicated in-app
administrator dashboard. The repository does contain administrative capabilities at the database
and deployment level (for example, \texttt{admin\_users} in SQL policies and Supabase admin APIs),
but operational administration is performed through deployment tooling, Supabase configuration, and
project infrastructure rather than through end-user screens. For that reason, this subsection
documents the administrator workflow that actually exists in the current system.

\subsubsection*{A. Initial platform setup}

\begin{enumerate}
  \item Create and configure the target Supabase project.
  \item Enable email/password auth and email confirmation.
  \item Enable Anonymous Auth if guest mode is required.
  \item Set the Site URL and redirect URLs for local and hosted environments.
  \item Apply the active database migrations in \texttt{supabase/migrations}.
  \item Deploy the two required Edge Functions:
  \begin{itemize}
    \item \texttt{material-worker}
    \item \texttt{guest-sandbox-cleanup}
  \end{itemize}
  \item Configure required secrets for the backend and both Edge Functions.
  \item Deploy the web app and Python backend.
\end{enumerate}

\subsubsection*{B. Key administrator responsibilities}

\begin{itemize}
  \item Maintain environment variables for the web app, backend, and Edge Functions.
  \item Apply the branded confirmation and recovery email templates in Supabase.
  \item Ensure \texttt{PYTHON\_BACKEND\_URL} and backend API-key alignment across services.
  \item Monitor guest-mode capacity and cleanup health.
  \item Run post-deploy smoke tests after releases.
  \item Investigate common failure modes such as:
  \begin{itemize}
    \item guest entry failures;
    \item materials stuck in \texttt{processing};
    \item auth links redirecting to the wrong domain;
    \item web UI loading while AI features fail due to backend or provider issues.
  \end{itemize}
\end{itemize}

\subsubsection*{C. Common admin commands}

From the current repository documentation:

\begin{itemize}
  \item Install dependencies:
  \begin{center}
  \texttt{pnpm install} \\
  \texttt{pip install -r backend/requirements.txt}
  \end{center}
  \item Run locally:
  \begin{center}
  \texttt{pnpm dev} \\
  \texttt{uvicorn app.main:app --app-dir backend --host 0.0.0.0 --port 8001 --reload}
  \end{center}
  \item Apply migrations:
  \begin{center}
  \texttt{npx supabase db push}
  \end{center}
\end{itemize}

\subsection{Instructor / Grader Manual}

\subsubsection*{A. Create an account and sign in}

\begin{enumerate}
  \item Open the landing page.
  \item Choose \textbf{Create an account}.
  \item Select the \textbf{teacher} account type.
  \item Enter email and password.
  \item Verify the email using the confirmation link.
  \item Sign in and open the \textbf{Teacher Dashboard}.
\end{enumerate}

\subsubsection*{B. Create and manage a class}

\begin{enumerate}
  \item From the Teacher Dashboard, click \textbf{Create class}.
  \item Fill in class title and optional metadata such as subject and level.
  \item Submit the form.
  \item The system creates the class and generates a join code.
  \item Share the join code with students so they can enroll from the student side.
\end{enumerate}

\subsubsection*{C. Upload teaching materials}

\begin{enumerate}
  \item Open the target class.
  \item Use the materials upload form to upload PDF, DOCX, or PPTX files.
  \item Wait for the system to process the material.
  \item Monitor status as \texttt{processing}, \texttt{ready}, or \texttt{failed}.
  \item When the material is ready, use the actions menu to preview, download, or delete it.
\end{enumerate}

\subsubsection*{D. Generate and publish the Course Blueprint}

\begin{enumerate}
  \item Open the class blueprint page.
  \item Trigger blueprint generation after materials are ready.
  \item Review the generated summary, topics, and objectives.
  \item Edit the blueprint draft where needed.
  \item Approve the draft and move it into overview.
  \item Publish the blueprint when it is ready for student use.
\end{enumerate}

\subsubsection*{E. Create learning activities}

The current release supports three main activity types:
\begin{itemize}
  \item chat assignments;
  \item quizzes;
  \item flashcards.
\end{itemize}

To create an activity:
\begin{enumerate}
  \item Ensure a published blueprint exists.
  \item Open the relevant activity creation route from the class page.
  \item Generate or configure the activity.
  \item Review and edit the content where applicable.
  \item Publish or assign it to the class.
\end{enumerate}

\subsubsection*{F. Review submissions}

\begin{enumerate}
  \item Open an assignment's review page.
  \item View student recipients, completion states, and submitted content.
  \item For chat assignments, inspect transcript and reflection content.
  \item For quiz assignments, inspect scores and answers.
  \item For flashcards, inspect review statistics and score summaries.
  \item Save teacher feedback for the selected submission.
\end{enumerate}

\subsubsection*{G. Monitor class progress and teaching support}

\begin{enumerate}
  \item Open the class overview to see the adaptive teaching brief.
  \item Refresh the brief when more activity data becomes available.
  \item Open \textbf{Insights} for the class intelligence dashboard.
  \item Review charts, student summaries, AI narrative, and intervention suggestions.
  \item Use the natural-language data query panel for custom chart-oriented questions.
\end{enumerate}

\subsubsection*{H. Preview the student experience}

\begin{enumerate}
  \item Open the class page as a teacher.
  \item Activate student preview mode using the student-view route state.
  \item Inspect how published blueprint, AI chat, and assignment widgets appear to students.
  \item Exit preview mode to return to the teacher control surface.
\end{enumerate}

\subsection{Student Manual}

\subsubsection*{A. Create an account and sign in}

\begin{enumerate}
  \item Open the landing page.
  \item Choose \textbf{Create an account}.
  \item Select the \textbf{student} account type.
  \item Enter email and password.
  \item Verify the email using the confirmation link.
  \item Sign in and open the \textbf{Student Dashboard}.
\end{enumerate}

\subsubsection*{B. Join a class}

\begin{enumerate}
  \item From the Student Dashboard or sidebar, open \textbf{Join class}.
  \item Enter the join code provided by the instructor.
  \item Submit the form.
  \item After joining successfully, open the class workspace.
\end{enumerate}

\subsubsection*{C. Use the student dashboard}

The student dashboard organizes assignments into categories such as:
\begin{itemize}
  \item due now;
  \item upcoming;
  \item completed.
\end{itemize}

Students can use the dashboard to quickly identify overdue or active work and jump directly into
the relevant class or assignment.

\subsubsection*{D. Navigate the class workspace}

Within a class, students can access:
\begin{itemize}
  \item AI chat;
  \item chat assignments;
  \item quizzes;
  \item flashcards;
  \item the published blueprint.
\end{itemize}

The workspace is intentionally more focused and less administratively dense than the teacher view.

\subsubsection*{E. Use AI chat}

\begin{enumerate}
  \item Open the class page and select the AI chat widget or direct chat view.
  \item Ask questions related to the class materials and published blueprint.
  \item Review the grounded AI response.
  \item Continue the conversation as needed.
\end{enumerate}

If chat is not yet available, the likely reason is that the instructor has not published the
blueprint yet.

\subsubsection*{F. Complete assignments}

\begin{enumerate}
  \item Open the assigned activity from the dashboard or class workspace.
  \item For chat assignments, complete the conversation and submit the required reflection.
  \item For quizzes, answer the questions and submit the attempt.
  \item For flashcards, complete the review session and submit progress.
  \item Revisit the assignment later to see updated status such as submitted or reviewed.
\end{enumerate}

\subsubsection*{G. Use the published blueprint}

\begin{enumerate}
  \item Open the blueprint widget or blueprint page inside the class.
  \item Review topics and learning objectives.
  \item Use the blueprint as a study guide before attempting quizzes or asking AI chat questions.
\end{enumerate}

\subsubsection*{H. Manage account settings}

\begin{enumerate}
  \item Open \textbf{Settings}.
  \item Update display name if needed.
  \item Use the two-step password change flow with verification code if required.
  \item Open \textbf{Help} for role-specific FAQ and suggested next steps.
\end{enumerate}

\subsection{Supplementary Guest Demo Flow}

Although the course report specifically asks for admin, instructor, and student manuals, the guest
workflow is worth documenting because it is a significant implemented feature.

\begin{enumerate}
  \item From the landing page, choose \textbf{Continue as guest} when guest mode is enabled.
  \item The system provisions an anonymous authenticated session and clones sandbox data.
  \item Explore the classroom in teacher view or switch to student view.
  \item Reset the sandbox if a fresh walkthrough is needed.
  \item Create a permanent account from a clean state if you want to move into normal usage.
\end{enumerate}

\section{Conclusion}

The current repository demonstrates a serious course project with a clear educational purpose,
a layered architecture, meaningful GenAI integration, and a credible software engineering process.
Its strongest design choice is the decision to center the product on a teacher-curated Course
Blueprint rather than on generic AI chat alone. That decision gives the entire system coherence:
materials, activities, student chat, review, and analytics all connect back to the same core object.

Equally important, the project goes beyond feature implementation and addresses the concerns that
often separate prototypes from robust systems: asynchronous processing, security boundaries,
role-aware access control, test coverage, deployment topology, and realistic demo-mode isolation.
For a final project, this makes the platform not only functional, but also defensible from a
software engineering perspective.

\appendix

\section{Repository Evidence Used for This Report}

This report was prepared against the current repository state using the following main sources:

\begin{itemize}
  \item top-level documents:
  \begin{itemize}
    \item \texttt{README.md}
    \item \texttt{ARCHITECTURE.md}
    \item \texttt{DESIGN.md}
    \item \texttt{UIUX.md}
    \item \texttt{DEPLOYMENT.md}
  \end{itemize}
  \item subsystem documentation:
  \begin{itemize}
    \item \texttt{web/README.md}
    \item \texttt{backend/README.md}
    \item \texttt{supabase/README.md}
  \end{itemize}
  \item course materials:
  \begin{itemize}
    \item \texttt{project\_overivew.pdf}
    \item \texttt{project\_background.pdf}
  \end{itemize}
  \item implementation evidence:
  \begin{itemize}
    \item key web routes and server actions
    \item backend service modules
    \item Supabase Edge Functions
    \item migration history
    \item git commit history and merge history
    \item test suites and E2E specs
  \end{itemize}
\end{itemize}

\section{Current Accuracy Caveat About the ``Admin'' Role}

The course deliverable asks for an admin user manual, but the current repository does not expose a
dedicated end-user admin console. Permanent in-app roles are teacher and student, while guest mode
is a sandboxed anonymous access path. Administrative responsibilities in the current release exist
mainly at the infrastructure and deployment layer. This report therefore documents the administrator
workflow that is actually implemented today, rather than inventing an unsupported UI role.

\end{document}
